Una cámara trampa es un sensor que dispara al detectar movimiento y guarda la
fotografía sin más contexto. Un proyecto de monitoreo despliega decenas y
acumula cientos de miles de imágenes por temporada, de las cuales la gran
mayoría no contiene ningún animal: una rama moviéndose con el viento o un cambio
de iluminación bastan para activar el disparador. Revisarlas a mano es el cuello
de botella real del monitoreo de fauna, y es la tarea que un detector puede
absorber.

El problema tiene tres dificultades que no son accidentales del conjunto elegido
sino propias del dominio, y las tres condicionan el diseño de este trabajo.

\emph{El 70\% de las imágenes está vacío.} Un detector no se entrena con
imágenes sin objeto, de modo que el conjunto útil es una fracción del total. Pero
la decisión de qué hacer con las vacías no es trivial: tratarlas como una clase
más sería un error conceptual, porque la ausencia de cajas ya expresa que no hay
nada que detectar.

\emph{El desbalance es de factor 500.} La especie más común tiene 2\,514 cajas
de entrenamiento y la más rara, cinco. Con cinco ejemplos no hay nada que
aprender, y una métrica agregada que promedie todas las especies por igual lo
dirá; una que pondere por frecuencia lo esconderá. Se reportan ambas.

\emph{El modelo tiende a memorizar el fondo.} Cada cámara está fija en un
poste y produce miles de imágenes del mismo paisaje. Un modelo que aprende
``cuando aparece algo en este tronco, es un zorro'' funciona en esa cámara y
falla en la siguiente. Beery et al.~\cite{beery2018recognition} documentaron el
fenómeno y diseñaron para medirlo la partición que este trabajo usa: conjuntos
de prueba en las \emph{mismas} ubicaciones del entrenamiento (\emph{cis}) y en
ubicaciones \emph{nunca vistas} (\emph{trans}).

Sobre esa base, el trabajo ajusta un detector YOLO y responde cinco preguntas
que el enunciado plantea: qué desempeño alcanza y cómo se reparte entre especies;
cuánto pierde al cambiar de ubicación; dónde mira realmente para decidir; cuánto
resiste la oclusión y la lejanía del animal; y qué reporte ecológico se puede
construir con sus salidas sin fingir una precisión que no tiene.

Las contribuciones son cuatro. Primero, una reimplementación propia de las
métricas que distingue
\emph{AP indefinida} de \emph{AP cero} para las especies ausentes de un conjunto
y recalcula cualquier subconjunto sin volver a correr el modelo. Segundo, un
control de distribución que separa la pérdida por cambio de fondo de la del
cambio de prevalencia de especies. Tercero, una adaptación de Grad-CAM a
detección con la capa elegida por medición y no por convención. Y cuarto, un
módulo de monitoreo que además de producir el reporte de actividad lo
\emph{valida} contra las anotaciones humanas, midiendo el error especie por
especie.
